\chapter{Urology}

\medterm{Bladder Carcinoma}-- Malignant neoplasm originating from the urothelium lining the urinary bladder, with transition cell carcinoma (urothelial carcinoma) accounting for over 90 percent of cases. Predominant risk factor: cigarette smoking (responsible for ~50\% of cases) and occupational exposure to aromatic amines or aniline dyes. Classic presentation: painless gross hematuria. Diagnosis: flexible cystoscopy with biopsy and urine cytology. Staging (TNM): non-muscle-invasive bladder cancer (Ta, T1, Tis) versus muscle-invasive bladder cancer ($\ge$T2). Treatment for non-muscle-invasive disease: transurethral resection of bladder tumor (TURBT) followed by intravesical Bacillus Calmette-Guérin (BCG) or chemotherapy. Treatment for muscle-invasive disease: radical cystectomy with urinary diversion (ileal conduit or neobladder) and neoadjuvant cisplatin-based combination chemotherapy.

\medterm{Bladder Health and Disorders}-- Bladder health involves the proper storage and voluntary release of urine through coordinated function of the detrusor muscle and sphincters. Common disorders include urinary tract infections, overactive bladder, stress incontinence, and interstitial cystitis. Maintaining bladder health includes adequate hydration, pelvic floor exercises, and prompt treatment of infections.

\medterm{Bladder Problems}-- a range of conditions affecting urine storage and voiding, including incontinence, overactive bladder, urinary retention, and recurrent infections. Causes may include neurological damage, pelvic floor weakness, prostate enlargement, or bladder outlet obstruction. Evaluation includes urinalysis, urodynamics, and imaging; treatment varies from behavioral modifications and medications to surgery.

\medterm{Catheterization (Urinary)}-- Insertion of a flexible tube (catheter) through the urethra or suprapubically into the urinary bladder to drain urine. Types: Foley catheter (indwelling balloon catheter), intermittent (in-out) catheter, and suprapubic catheter. Indications: acute urinary retention, monitoring hourly urine output in critically ill patients, intraoperative management, neurogenic bladder, or intractable incontinence when conservative measures fail. Complications: catheter-associated urinary tract infection (CAUTI), urethral trauma, stricture formation, bladder spasms, and balloon malfunction. Prevention of CAUTI requires strict aseptic insertion technique, maintaining a closed drainage system, and prompt removal when no longer clinically indicated.

\medterm{Circumcision}-- Surgical removal of the foreskin (prepuce) covering the glans penis. Indications: medical reasons such as severe phimosis, recurrent balanoposthitis, paraphimosis, or lichen sclerosus (balanitis xerotica obliterans); elective infant circumcision for religious, cultural, or preventive health reasons. Preventive benefits: reduced risk of urinary tract infections in male infants, reduced risk of heterosexually acquired HIV transmission (demonstrated in sub-Saharan African clinical trials), decreased incidence of penile cancer, and lower risk of human papillomavirus (HPV) and herpes simplex virus type 2 transmission to partners. Surgical methods: Gomco clamp, Plastibell device, Mogen clamp, or dorsal slit and sleeve resection in older males.

\medterm{Cystitis}-- Inflammation of the urinary bladder, most frequently caused by bacterial infection originating from colonic flora. Escherichia coli is responsible for 75--90\% of uncomplicated cases, followed by Staphylococcus saprophyticus, Klebsiella pneumoniae, and Proteus mirabilis. More common in females due to shorter urethral length and proximity of the urethral meatus to the anus. Clinical features: dysuria, urinary frequency, urgency, suprapubic pain, and hematuria. Urinalysis shows pyuria, positive leukocyte esterase, and positive nitrites. Uncomplicated cystitis in non-pregnant women is treated with nitrofurantoin, trimethoprim-sulfamethoxazole, or fosfomycin. Non-infectious cystitis causes include radiation therapy, cyclophosphamide chemotherapy, and interstitial cystitis (painful bladder syndrome).

\medterm{Cystoscopy}-- Endoscopic examination of the lower urinary tract, allowing direct visualization of the urethra, prostate, bladder neck, ureteral orifices, and bladder mucosa. Equipment: flexible or rigid cystoscope with fiberoptic light source, lens system, and working channels for instruments. Flexible cystoscopy is performed under local anesthesia (topical lidocaine jelly) for diagnostic evaluation of hematuria, recurrent UTIs, bladder pain, or voiding dysfunction. Rigid cystoscopy is performed under regional or general anesthesia for operative procedures including bladder tumor resection (TURBT), ureteral stent placement/extraction, bladder stone lithotripsy, and mucosal biopsy.

\medterm{Dysuria}-- Painful, burning, or uncomfortable sensation experienced during urination. Most commonly indicates inflammation or infection of the lower urinary tract. Causes in females: acute cystitis, urethritis (Chlamydia trachomatis, Neisseria gonorrhoeae), vaginitis (Candida, Trichomonas), and vulvar atrophy. Causes in males: urethritis, prostatitis, benign prostatic hyperplasia, or epididymitis. Non-infectious causes: bladder calculi, urethral trauma, chemical irritation (soaps, spermicides), interstitial cystitis, and urothelial neoplasms. Initial diagnostic evaluation involves urinalysis, urine culture, and testing for sexually transmitted infections when indicated.

\medterm{Ejaculation}-- The expulsion of seminal fluid from the male urethra during sexual climax, a reflex coordinated by the sympathetic nervous system (T10--L2) and involving contraction of the epididymis, vas deferens, seminal vesicles, prostate, and bulbospongiosus muscle. The process has two phases: emission (deposition of sperm and seminal fluid into the posterior urethra) and expulsion (rhythmic contraction of the pelvic floor muscles propelling semen through the urethra). Ejaculatory disorders include premature ejaculation (most common), delayed ejaculation, retrograde ejaculation (semen flows backward into the bladder), anejaculation (no ejaculation), and painful ejaculation. Causes: neurologic, pharmacologic (SSRIs, alpha-blockers), psychologic, and structural.

\medterm{Emission}-- The discharge or release of a substance, usually a fluid.

\medterm{Emphysematous Pyelonephritis}-- Severe necrotizing infection of renal parenchyma with gas formation. Almost exclusively in diabetic patients. Presents with high fever, flank pain, and sepsis. CT shows gas within renal parenchyma. A urological emergency requiring aggressive antibiotics and usually nephrectomy. Mortality is high even with treatment.

\medterm{Epididymitis}-- Inflammation of the epididymis, a common cause of acute scrotal pain in adults, typically resulting from infection. In sexually active men under 35 years, the most common pathogens are sexually transmitted: Chlamydia trachomatis (most common) and Neisseria gonorrhoeae, with epididymitis often accompanied by urethritis. In men over 35 years and in children, the causative organisms are enteric gram-negative bacilli (Escherichia coli, Pseudomonas, Klebsiella) from urinary tract infection, prostatitis, or following urological instrumentation. Presents with gradual onset of scrotal pain and swelling, often with fever and dysuria. Diagnosis is clinical, supported by urine analysis and NAAT testing for sexually transmitted infections. Treatment includes appropriate antibiotics targeting the causative organism, scrotal elevation, and NSAIDs for symptom control.

\medterm{Erectile Dysfunction}-- The persistent inability to achieve or maintain an erection sufficient for satisfactory sexual performance. Affects up to 50\% of men aged 40-70. Etiology: organic (vascular disease, diabetes, hypogonadism, neurological disorders, medications, Peyronie disease) and psychogenic (anxiety, depression, relationship issues). Vascular causes: impaired arterial inflow (atherosclerosis) or venous leak. Diagnosis: history, IIEF-5 questionnaire, nocturnal penile tumescence testing, duplex Doppler ultrasound, and hormonal profile. Treatment: PDE5 inhibitors (sildenafil, tadalafil, vardenafil) as first-line; second-line options: intracavernosal injections (alprostadil), vacuum erection devices, and penile prosthesis for refractory cases. Lifestyle modification (weight loss, exercise, smoking cessation) improves outcomes.

\medterm{Erectile Dysfunction (ED)}-- Inability to attain or maintain a penile erection sufficient for satisfactory sexual performance. Etiologies: vasculogenic (atherosclerosis, hypertension, diabetes mellitus), neurogenic (radical prostatectomy, spinal cord injury, multiple sclerosis, diabetic neuropathy), hormonal (hypogonadism, hyperprolactinemia), psychogenic (performance anxiety, depression, relationship conflict), and drug-induced (antihypertensives, antidepressants, antiandrogens). Vascular ED serves as an early marker for systemic coronary artery disease due to endothelial dysfunction in smaller cavernosal arteries. First-line therapy: oral phosphodiesterase type 5 (PDE5) inhibitors (sildenafil, tadalafil, vardenafil). Second-line therapy: intracavernosal alprostadil injection or vacuum erection devices. Third-line therapy: surgical implantation of a inflatable penile prosthesis.

\medterm{Extracorporeal Shock Wave Lithotripsy (ESWL)}-- Non-invasive treatment method for urinary calculi that uses focused acoustic shock waves generated outside the body to fragment kidney or ureteral stones into smaller gravel-like pieces (dust/sand) that pass spontaneously in urine. Shock wave sources: electrohydraulic, electromagnetic, or piezoelectric generators, focused using fluoroscopic or ultrasound guidance. Most effective for renal stones $<2$~cm and proximal ureteral stones with density $<1000$~Hounsfield Units on non-contrast CT. Contraindications: pregnancy, uncorrected bleeding diatheses, uncontrolled urinary tract infection, severe body habitus, and abdominal aortic aneurysm. Complications: renal hematoma, steinstrasse (ureteral obstruction by stone fragments), and transient hematuria.

\medterm{Hematuria}-- Presence of red blood cells in the urine, categorized as microscopic ($>3$~RBCs per high-power field on urine microscopy) or gross (macroscopic, visible color change from pink to dark red). Microscopic hematuria is often asymptomatic and detected on routine urinalysis. Gross hematuria requires urgent urological investigation to exclude malignancy. Etiologies: urological sources (bladder carcinoma, renal cell carcinoma, prostate cancer, nephrolithiasis, UTI, BPH, trauma) versus nephrological sources (IgA nephropathy, post-infectious glomerulonephritis, Alport syndrome). Glomerular origin is suggested by dysmorphic RBCs, red cell casts, and significant proteinuria. Workup for non-glomerular hematuria: multiphase CT urography and flexible cystoscopy.

\medterm{Hydrocele}-- Accumulation of serous fluid within the tunica vaginalis surrounding the testis, presenting as a painless, soft, fluid-filled scrotal swelling. Types: congenital (communicating hydrocele due to a patent processus vaginalis, common in infants, fluctuates in size) and acquired (non-communicating, resulting from impaired fluid absorption due to epididymitis, orchitis, testicular trauma, filariasis, or tumor). Diagnosis: scrotal transillumination (hydrocele glows brightly red) and scrotal ultrasonography to rule out underlying testicular neoplasm. Treatment: conservative observation for asymptomatic cases; hydrocelectomy (Jaboulay or Lord technique) or aspiration with sclerotherapy for large, symptomatic, or tense hydroceles.

\medterm{Hypospadias}-- Congenital anomaly of male urogenital development characterized by an abnormally proximal urethral meatus located on the ventral surface of the penis, scrotum, or perineum, rather than at the tip of the glans. Classified by meatal location: anterior (glandular, subcoronal—60--70\%), middle (penile), or posterior (equinoscrotal, perineal). Associated features: ventral curvature of the penis (chordee) and an incomplete, hooded foreskin. Resulting from failure of complete fusion of the urethral folds during embryogenesis (weeks 8--14). Routine infant circumcision is contraindicated because the prepuce is required for surgical reconstruction (hypospadias repair/urethroplasty), typically performed between 6 and 18 months of age.

\medterm{Kidney Stones (Nephrolithiasis / Urolithiasis)}-- Crystalline mineral deposits formed within the renal pelvis and calyces due to supersaturation of urine. Types: calcium oxalate/phosphate (75--80\%, radio-opaque), struvite (magnesium ammonium phosphate, 10--15\%, staghorn calculi associated with Proteus infection), uric acid (5--10\%, radiolucent on X-ray, associated with gout and metabolic syndrome), and cystine (1--2\%, ground-glass appearance, autosomal recessive cystinuria). Symptoms: sudden, severe renal colic radiating from flank to groin, hematuria, nausea, vomiting, and dysuria. Diagnosis: non-contrast low-dose abdominal/pelvic CT (gold standard). Management: conservative passage with hydration and alpha-blockers for stones $<5$~mm; ESWL, ureteroscopy with laser lithotripsy, or percutaneous nephrolithotomy (PCNL) for larger or non-passing stones.

\medterm{Libido and Sexual Performance}-- Libido (sexual desire) is modulated by a complex interplay of hormonal (testosterone, estrogen, prolactin), neurologic, psychologic, and relational factors. Low libido in men is commonly associated with hypogonadism, depression, medication side effects (SSRIs, antihypertensives), and chronic illness; in women, contributing factors include hormonal changes (menopause, postpartum), relationship distress, stress, and body image concerns. Sexual performance encompasses erectile function, ejaculatory control, and orgasmic capacity. Evaluation includes validated questionnaires, laboratory testing (testosterone, prolactin, thyroid function, glucose, lipids), and psychologic assessment. Treatment targets the underlying cause: testosterone replacement for hypogonadism, PDE5 inhibitors for ED, SSRIs or topical anesthetics for premature ejaculation, and psychosexual therapy or couples counseling for relationship or psychogenic contributors.

\medterm{Lithotripsy}-- A medical procedure that uses shock waves to fragment kidney stones (renal calculi) into smaller pieces that can pass through the urinary tract. See Shock Wave Lithotripsy.

\medterm{Masturbation}-- Self-stimulation of the genitalia for sexual pleasure, a normal and common human behavior across all ages and genders. Prevalence: >90\% of males and >60\% of females report having masturbated at some point in their lives. It is a healthy component of psychosexual development, stress relief, and sexual self-exploration. There are no adverse medical consequences. Excessive or compulsive masturbation may rarely be associated with underlying psychiatric conditions (e.g., obsessive-compulsive disorder, hypersexuality). Cultural and religious attitudes vary widely. Medical benefits: may reduce stress, improve sleep, relieve sexual tension, and (in males) may reduce the risk of prostate cancer by facilitating the clearance of prostatic secretions.

\medterm{Nephrolithiasis}-- Stone formation in the urinary tract. Calcium oxalate most common at 80\%, followed by struvite, uric acid, and cystine. Risk factors include dehydration, hypercalciuria, hyperoxaluria, hypocitraturia, and low urine volume. Evaluation includes 24-hour urine analysis. Treatment includes hydration, dietary modification, and targeted medical therapy based on stone composition.

\medterm{Nocturia}-- Complaint of waking one or more times during the night to void, with each void preceded and followed by sleep. Highly prevalent in older adults and significantly impairs sleep quality and quality of life. Etiologies: nocturnal polyuria (nocturnal urine volume $>33\%$ of total 24-hour volume, caused by heart failure, obstructive sleep apnea, peripheral edema, or loss of circadian vasopressin secretion), low nocturnal bladder capacity (BPH, neurogenic bladder, interstitial cystitis, overactive bladder), or global polyuria (diabetes mellitus, diabetes insipidus, primary polydipsia). Evaluation: 24-hour frequency-volume chart (bladder diary). Treatment targets underlying systemic or urological causes; desmopressin may be used for nocturnal polyuria in selected patients.

\medterm{Orchitis}-- Inflammation of one or both testes, most commonly resulting from hematogenous viral infection or contiguous extension from acute epididymitis (epididymo-orchitis). Viral orchitis: classic complication of mumps (paramyxovirus), occurring in up to 30\% of post-pubertal males 4--8 days after parotitis. Bacterial orchitis: typically caused by Chlamydia trachomatis or Neisseria gonorrhoeae in sexually active men under 35, and enteric organisms (E. coli, Pseudomonas) in older men. Clinical features: acute scrotal pain, testicular swelling, fever, systemic toxicity, and reactive hydrocele. Complications: testicular atrophy, impaired spermatogenesis, and infertility. Treatment: supportive care (analgesics, bed rest, scrotal elevation) for viral orchitis; empirical antibiotics for bacterial causes.

\medterm{Paraphimosis}-- Urological emergency in uncircumcised males occurring when a retracted, tight foreskin becomes trapped behind the coronal sulcus of the glans penis, forming a constricting ring. Venous and lymphatic congestion causes progressive edema of the glans and prepuce, eventually impairing arterial perfusion and risking glans ischemia, necrosis, and gangrene. Causes: failure to return the foreskin to its anatomical position following urinary catheterization, hygiene, or physical examination. Management: immediate manual reduction using digital compression to reduce edema followed by traction on the prepuce; if manual reduction fails, ice packs, osmotic agents (granulated sugar), hyaluronidase injections, dorsal slit procedure, or emergency circumcision are indicated.

\medterm{Percutaneous Nephrolithotomy (PCNL)}-- Minimally invasive surgical procedure for removing large, complex, or staghorn renal calculi. Access: a needle track is established under fluoroscopic or ultrasound guidance from the flank through the renal parenchyma into the target renal calyx, dilated to allow passage of a nephroscope and working sheath. Stones are fragmented using ultrasonic, pneumatic, or laser energy sources and extracted under direct vision. Indications: kidney stones $>2$~cm, staghorn calculi, lower pole stones $>1.5$~cm, or cases where ESWL and ureteroscopy have failed. Complications: hemorrhage requiring transfusion, retroperitoneal hematoma, pleural entry (hydro/pneumothorax), sepsis, and adjacent organ injury.

\medterm{Phimosis}-- Inability to retract the foreskin (prepuce) proximally over the glans penis. Physiologic phimosis is normal in infants and uncircumcised young boys due to natural adhesions between glans and prepuce, resolving spontaneously in over 90\% by age 3--5 years. Pathologic phimosis is secondary to scarring, inflammation, or infection, presenting in adolescents or adults as a tight, fibrotic preputial ring. Causes of pathologic phimosis: recurrent balanoposthitis, poor hygiene, severe diabetes, and lichen sclerosus (balanitis xerotica obliterans). Clinical features: painful erections, ballooning of foreskin during micturition, recurrent UTIs, and increased risk of penile cancer. Treatment: topical high-potency corticosteroid cream (betamethasone) with manual stretching; surgical circumcision or preputioplasty for refractory cases.

\medterm{Priapism}-- Prolonged, painful erection lasting longer than 4 hours, occurring independently of sexual desire or stimulation. Medical emergency categorized into two types: ischemic (low-flow, veno-occlusive) and non-ischemic (high-flow, arterial). Ischemic priapism accounts for $>95\%$ of cases, characterized by compartment syndrome of the corpora cavernosa with stasis, hypoxia, and acidosis; causes include sickle cell disease, intracavernosal erectile dysfunction therapies, medications (trazodone, antipsychotics), and hematological malignancies. Non-ischemic priapism is usually secondary to perineal/penile trauma resulting in an arteriocavernosal fistula. Management of ischemic priapism: cavernous blood aspiration, intracavernosal injection of alpha-adrenergic agonists (phenylephrine), and surgical shunting procedures (Winter, Al-Ghorab, or distal shunt) if conservative measures fail after 12--24 hours to prevent irreversible erectile tissue fibrosis and impotence.

\medterm{Prostatectomy}-- Surgical procedure involving the partial or total removal of the prostate gland. Simple prostatectomy: removal of the inner adenoma while leaving the peripheral capsule intact, performed open, laparoscopically, or robotically for very large BPH ($>80$--100~grams). Radical prostatectomy: complete excision of the entire prostate gland, seminal vesicles, and ampullae of the vas deferens, along with pelvic lymph node dissection for localized prostate cancer. Radical prostatectomy approaches: open retropubic, perineal, or robot-assisted laparoscopic radical prostatectomy (RALP). Key surgical complications: urinary stress incontinence (due to external urethral sphincter trauma) and erectile dysfunction (due to cavernous nerve injury; nerve-sparing techniques preserve potency when oncologically feasible).

\medterm{Prostate-Specific Antigen (PSA)}-- Glycoprotein enzyme (serine protease) produced almost exclusively by the luminal epithelial cells of the prostate gland, secreted into seminal fluid to liquefy the semen coagulum. Small amounts enter the bloodstream, where total PSA can be measured. Reference range: generally $<4$~ng/mL, though age-adjusted thresholds are used. Elevated PSA occurs in prostate cancer, but is non-specific and also elevated in benign prostatic hyperplasia, acute/chronic prostatitis, urinary retention, prostate biopsy, and recent ejaculation. Clinical utility: screening (discussed with informed decision-making), monitoring response to prostate cancer treatment, and surveillance for disease recurrence. Diagnostic specificity is improved using free-to-total PSA ratio, PSA velocity, and PSA density.

\medterm{Prostatitis}-- Inflammation or infection of the prostate gland, classified by the National Institutes of Health (NIH) into four categories: Category I (Acute Bacterial Prostatitis), Category II (Chronic Bacterial Prostatitis), Category III (Chronic Prostatitis / Chronic Pelvic Pain Syndrome, CP/CPPS), and Category IV (Asymptomatic Inflammatory Prostatitis). Acute bacterial prostatitis presents with fever, chills, perineal/low back pain, dysuria, and acute voiding dysfunction; prostate is tender, warm, and boggy on DRE (prostate massage is contraindicated to avoid bacteremia). Caused by Gram-negative bacilli (E. coli, Pseudomonas). Category III (CP/CPPS) accounts for 90\% of symptomatic cases, treated with multimodal therapy (alpha-blockers, anti-inflammatories, physical therapy).

\medterm{Pyelonephritis}-- Acute or chronic infection of the renal parenchyma and pelvicalyceal system, most commonly resulting from ascending urinary tract infection from the bladder. Major pathogens: Escherichia coli ($>80\%$), Proteus, Klebsiella, Enterobacter, and Pseudomonas. Risk factors: female gender, pregnancy, urinary tract obstruction (stones, BPH), vesicoureteral reflux, neurogenic bladder, and diabetes mellitus. Clinical features: triad of flank pain/CVA tenderness, high fever with rigors, and lower urinary tract symptoms (dysuria, frequency). Urinalysis shows WBC casts (pathognomonic for upper UTI), pyuria, and bacteriuria. Complications: renal abscess, emphysematous pyelonephritis, sepsis, and chronic renal scarring. Treatment: empiric systemic antibiotic therapy (fluoroquinolones, third-generation cephalosporins, or aminoglycosides).

\medterm{Renal Cell Carcinoma (RCC)}-- Primary malignant adenocarcinoma of the renal cortex, originating from proximal renal tubular epithelial cells. Represents ~85\% of all primary adult kidney cancers. Major histologic subtypes: clear cell RCC (70--80\%, associated with VHL gene mutations on chromosome 3p), papillary RCC (10--15\%), and chromophobe RCC (5\%). Risk factors: cigarette smoking, obesity, hypertension, end-stage renal disease, and occupational exposure to cadmium or asbestos. Classic triad ($<10\%$ of patients): flank pain, palpable abdominal mass, and gross hematuria; most cases are discovered incidentally on imaging. Paraneoplastic syndromes: erythrocytosis (EPO hypersecretion), hypercalcemia (PTHrP), hypertension (renin), and Stauffer syndrome (non-metastatic hepatic dysfunction). Treatment: partial or radical nephrectomy; metastatic RCC responds to targeted anti-angiogenic agents (tyrosine kinase inhibitors like sunitinib) and immune checkpoint inhibitors (nivolumab, ipilimumab).

\medterm{Semen}-- The fluid ejaculated from the male reproductive tract, containing sperm (produced in the testes) suspended in seminal plasma (produced by the seminal vesicles, prostate, and bulbourethral glands). Semen provides nutrients (fructose, amino acids), protective agents (zinc, antioxidants), and a medium for sperm transport. Normal semen parameters (WHO 2021): volume \(\ge\)1.4 mL, sperm concentration \(\ge\)16 million/mL, total sperm count \(\ge\)39 million, progressive motility \(\ge\)30\%, normal morphology \(\ge\)4\%.

\medterm{Shock Wave Lithotripsy}-- Non-invasive procedure using extracorporeal shock waves to fragment renal and ureteral stones. The patient is positioned over a water cushion coupling the shock wave generator to the body. Shock waves are focused on the stone under fluoroscopic or ultrasound guidance. Most effective for stones less than two centimeters in the renal pelvis or upper ureter. Success rates decrease with larger stones and lower pole location. Complications include renal hematoma, stone street from multiple fragments, residual stones, and need for multiple sessions, though the non-invasive nature makes it the preferred first-line treatment for suitable stones.

\medterm{Sperm}-- The male gamete (reproductive cell), produced in the seminiferous tubules of the testis through spermatogenesis. Sperm consists of: head (contains the haploid nucleus and acrosome containing enzymes for egg penetration), midpiece (contains mitochondria for energy), and tail (flagellum for motility). Approximately 1,500 sperm are produced per second. Sperm travel from the testis through the epididymis (maturation), vas deferens, seminal vesicles, and urethra during ejaculation.

\medterm{Spermatocele}-- Benign, painless, fluid-filled cystic retention cavity arising from the epididymis (usually the head/caput), containing avascular fluid and non-viable spermatozoa. Etiology: obstruction of an efferent ductule of the epididymis. Clinical presentation: smooth, spherical, non-tender mass located superior and posterior to the testis, distinct from the testicular parenchyma. Transilluminates clearly on physical examination. Scrotal ultrasound confirms a well-demarcated anechoic lesion separate from the testicle. Management: reassurance and conservative observation; surgical spermatocelectomy is reserved for large, painful, or bothersome cysts.

\medterm{Testicular Disorders}-- conditions affecting the testes, including testicular torsion (spermatic cord twisting causing ischemic pain and swelling, a surgical emergency requiring detorsion within 6 hours for testicular salvage), epididymitis (inflammation of the epididymis, often from bacterial infection—STI in young men, urinary pathogens in older men—treated with antibiotics and supportive care), orchitis (testicular inflammation, often viral from mumps), hydrocele (fluid collection in the tunica vaginalis, benign), varicocele (dilated pampiniform plexus veins, the most common correctable cause of male infertility), and testicular cancer (most common solid malignancy in young men 15–35, typically presenting as a painless testicular mass, diagnosed with scrotal ultrasound and tumor markers—AFP, hCG, LDH—and treated with radical inguinal orchiectomy with possible chemotherapy or radiation). Evaluation includes scrotal ultrasound, urinalysis, STI screening, and tumor markers.

\medterm{Testicular Torsion}-- Surgical emergency caused by twisting of the spermatic cord around its longitudinal axis, cutting off venous return and leading to arterial occlusion, ischemia, and infarction of the testicle. Predisposing anatomy: "bell-clapper" deformity, where the tunica vaginalis completely encircles the epididymis and distal cord, allowing the testis to rotate freely within the scrotum. Most common in neonates and post-pubertal males (ages 12--18). Symptoms: sudden onset of severe, unilateral scrotal pain, nausea, vomiting, high-riding testis with horizontal lie, and absent cremasteric reflex. Negative Prehn sign (pain is not relieved by scrotal elevation). Color Doppler ultrasound shows absent or decreased intratesticular blood flow. Management: urgent surgical exploration with manual detorsion and bilateral orchidopexy; testicular salvage rate is $>90\%$ if operated within 6 hours, dropping significantly after 12--24 hours.

\medterm{Transurethral Resection of the Prostate (TURP)}-- Gold standard surgical procedure for benign prostatic hyperplasia (BPH) refractory to medical management or complicated by urinary retention, renal insufficiency, or recurrent hematuria. Performed using a rigid resectoscope inserted through the urethra; a high-frequency electrosurgical wire loop (monopolar or bipolar) is used to systematically resect obstructive hyperplastic prostatic tissue from the transition zone under direct vision. Resected tissue chips are irrigated out and sent for histopathology. Complications: bleeding, retrograde ejaculation ($>75\%$), urethral stricture, urinary incontinence, and TURP syndrome (hyponatremia, fluid overload, and neurological symptoms resulting from systemic absorption of hypotonic monopolar irrigation fluid like glycine).

\medterm{Undescended Testis (Cryptorchidism)}-- Congenital failure of one or both testes to descend along the normal path from the retroperitoneum into the scrotum during fetal development. Most common congenital anomaly of the male genitalia, affecting 3\% of full-term and up to 30\% of premature male infants. Spontaneous descent usually occurs within the first 3--4 months of life. If uncorrected, cryptorchidism leads to impaired spermatogenesis, testicular atrophy, male infertility, increased risk of inguinal hernia, testicular torsion, and a 5- to 10-fold increased risk of testicular germ cell tumors (seminoma). Treatment: surgical orchidopexy (repositioning and anchoring the testis in the scrotal pouch), ideally performed between 6 and 18 months of age.

\medterm{Ureteral Obstruction}-- Blockage of urine flow causing proximal dilation and increased pressure. Acute causes renal colic with flank pain. Chronic leads to hydronephrosis and CKD. Emergent decompression with stent or nephrostomy for infected obstructed kidney. Diagnosis by CT or ultrasound showing hydronephrosis.

\medterm{Ureteral Stent (Double-J Stent)}-- A tube placed in the ureter from the renal pelvis to the bladder to maintain patency, with pigtail coils at both ends (hence double-J). Indications: ureteral obstruction (stones, stricture, malignancy), after ureteroscopy (prevent post-op obstruction), and as a pre-treatment for complex ureteral stones. Placed cystoscopically or antegrade through a nephrostomy. Sizes: 6-8 Fr, lengths 22-30 cm. Left in for weeks to months (though should not exceed 3 months due to encrustation). Complications: stent-related discomfort (urgency, frequency, hematuria, flank pain — very common), infection, migration, encrustation (forgotten stent), and stone formation.

\medterm{Ureteroscopy}-- Endoscopic procedure utilizing a thin, flexible or semirigid ureteroscope passed through the urethra, bladder, and into the ureter or renal collecting system. Used primarily for the diagnostic evaluation and therapeutic management of ureteral and renal calculi, ureteral strictures, and upper tract urothelial carcinoma (UTUC). Therapeutic ureteroscopy incorporates laser fiber lithotripsy (holmium:YAG or thulium fiber laser) to pulverize stones into dust, followed by stone basket retrieval. A temporary ureteral stent (Double-J stent) is frequently placed postoperatively to ensure urine drainage and prevent ureteral spasm or stricture formation.

\medterm{Urethral Stricture}-- Narrowing of the urethra due to fibrosis of the corpus spongiosum, most commonly caused by trauma (pelvic fracture, straddle injury, iatrogenic from catheterisation or urethral instrumentation), infection (gonococcal urethritis leading to stricture years later, now less common in the antibiotic era), and lichen sclerosus (Balanitis Xerotica Obliterans, BXO). Presents with progressive urinary obstruction: weak stream, hesitancy, straining, incomplete bladder emptying, frequency, and recurrent UTIs. Complete obstruction causes acute urinary retention. Diagnosis: retrograde urethrogram demonstrates the location and length of the stricture. Urethroscopy confirms and may allow simultaneous treatment. Treatment depends on location, length, and severity: urethral dilation (temporary, high recurrence rate), internal urethrotomy (optical incision of the stricture, 50\% recurrence at 2 years for bulbar strictures <2 cm), and urethroplasty (open surgical reconstruction, the gold standard for recurrent or long strictures, success rates >90\% for primary bulbar repairs). Complications include recurrence, erectile dysfunction, and incontinence.

\medterm{Urethritis}-- Inflammation of the urethra, presenting with dysuria, urethral pruritus, and urethral discharge. Classified into Gonococcal Urethritis (GU), caused by Neisseria gonorrhoeae, and Non-Gonococcal Urethritis (NGU), most frequently caused by Chlamydia trachomatis, followed by Mycoplasma genitalium, Trichomonas vaginalis, and Ureaplasma urealyticum. Transmission is primarily via sexual contact. Diagnosis: Gram stain or methylene blue stain of urethral smear showing $\ge$2--5 PMNs per high-power field, paired with nucleic acid amplification testing (NAAT) for N meisseria and Chlamydia. Treatment: empirical dual therapy covering both pathogens (e.g., ceftriaxone IM plus oral doxycycline or azithromycin), along with partner notification and treatment.

\medterm{Urinary Bladder}-- See Bladder Health and Disorders / Bladder Problems.

\medterm{Urinary Incontinence}-- Involuntary leakage of urine, significantly affecting psychological, social, and physical well-being. Main types: Stress Urinary Incontinence (SUI—leakage with increased intra-abdominal pressure such as coughing, sneezing, laughing; caused by urethral hypermobility or intrinsic sphincter deficiency), Urge Urinary Incontinence (UUI—involuntary leakage accompanied by a sudden, intense desire to void; caused by detrusor overactivity), Mixed Urinary Incontinence (combination of SUI and UUI), and Overflow Urinary Incontinence (involuntary dribbling due to urinary retention and bladder overdistension; caused by BPH, urethral stricture, or neurogenic detrusor underactivity). Management: pelvic floor muscle training (Kegel exercises), lifestyle modifications, anticholinergics/beta-3 agonists (mirabegron) for UUI, and surgical slings (mid-urethral sling) or artificial urinary sphincter for SUI.

\medterm{Urinary Retention}-- Inability to completely empty the bladder, which can be acute (sudden, painful, inability to void) or chronic (gradual, often painless, with overflow incontinence). Causes: benign prostatic hyperplasia, urethral stricture, neurogenic bladder, medications (anticholinergics), and pelvic organ prolapse. Diagnosis: post-void residual >100 mL, ultrasound, cystoscopy. Acute treatment: urethral catheterization. Chronic: alpha-blockers, 5-alpha-reductase inhibitors, and surgery.

\medterm{Urinary Tract Infection (UTI)}-- Inflammatory response of the urothelium to microbial invasion, ranging from lower tract infection (cystitis, urethritis) to upper tract infection (pyelonephritis, renal abscess). Classification: uncomplicated (infection in a healthy, non-pregnant, immunocompetent adult female with normal urinary anatomy) versus complicated (infection associated with factors increasing treatment failure or tissue damage, such as male gender, pregnancy, urinary obstruction, neurogenic bladder, indwelling catheter, immunosuppression, or renal insufficiency). Escherichia coli causes $>80\%$ of all community-acquired UTIs. Clinical assessment, urinalysis, urine culture, and targeted antimicrobial therapy form the core management strategy.

\medterm{Urodynamic Testing}-- Suite of diagnostic procedures designed to evaluate the physiological function of the lower urinary tract (bladder and urethral sphincter) during filling, storage, and emptying phases. Components: Uroflowmetry (measures urine flow rate and volume), Cystometrogram (CMG—evaluates bladder compliance, sensation, capacity, and involuntary detrusor contractions during saline filling), Pressure-Flow Studies (differentiates bladder outlet obstruction from detrusor muscle underactivity during voiding), Urethral Pressure Profiling, and Electromyography (EMG—assesses sphincter muscle activity and detrusor-sphincter dyssynergia). Indications: neurogenic bladder, refractory urinary incontinence, failed BPH surgery, and complex voiding dysfunction.

\medterm{Urolithiasis}-- Stone formation in the urinary tract (kidneys, ureters, bladder, urethra). Types: calcium oxalate (most common, 60-70\%), calcium phosphate (15-20\%), struvite (infection-related, 10-15\%), uric acid (5-10\%), and cystine (1-2\%). Risk factors: dehydration, hypercalciuria, hyperoxaluria, hyperuricosuria, hypocitraturia, and urinary tract infections. Ureteral stones present with colicky flank pain radiating to groin, hematuria, nausea/vomiting. Diagnosis: non-contrast CT (gold standard), ultrasound. Treatment: medical expulsive therapy (tamsulosin) for distal <10mm stones; ESWL, ureteroscopy with laser lithotripsy, or PCNL for larger/complex stones.

\medterm{Urology}-- Surgical and medical specialty concerned with the study, diagnosis, and management of diseases affecting the male and female urinary tract (kidneys, ureters, urinary bladder, urethra) and the male reproductive system (prostate, testes, epididymis, seminal vesicles, penis). Subspecialties: urologic oncology, endourology and stone disease, female pelvic medicine and reconstructive surgery (FPMRS), neurourology and urodynamics, andrology and male infertility, pediatric urology, and renal transplantation. Common diagnostic modalities include urinalysis, cystoscopy, ultrasound, CT urography, and urodynamics.

\medterm{Varicocele}-- Abnormal dilation and tortuosity of the pampiniform venous plexus within the spermatic cord, resulting from incompetent or absent venous valves. Occurs in 15\% of adolescent and adult males; overwhelmingly common on the left side ($>85\%$) due to the left testicular vein entering the left renal vein at a right angle (90 degrees) under higher hydrostatic pressure, compared to the right testicular vein entering the inferior vena cava obliquely. Described on physical examination as feeling like a "bag of worms", which decompresses when supine and accentuates during a Valsalva maneuver. Leading treatable cause of male infertility (present in 40\% of infertile men), causing hyperthermia, venous stasis, and oxidative stress that impairs spermatogenesis. Treatment: microscopic varicocelectomy or percutaneous embolization for symptomatic pain, testicular atrophy, or abnormal semen parameters.

\medterm{Vasectomy}-- Minor surgical procedure for permanent male contraception involving surgical interruption, occlusion, or excision of a segment of both vasa deferentia, preventing spermatozoa from entering the ejaculate. Performed under local anesthesia using traditional incisional or no-scalpel vasectomy (NSV) techniques. Vas occlusion methods: cauterization, ligation with clips/sutures, or fascial interposition. Contraceptive reliance is not immediate; patients must use alternative contraception until a post-vasectomy semen analysis (PVSA) confirms azoospermia or $<100,000$ non-motile sperm/mL, typically evaluated 8--12 weeks and 20 ejaculations post-procedure. Reversal (vasovasostomy) is possible but technically challenging and not guaranteed to restore fertility.

\medterm{Vesicoureteral Reflux}-- Retrograde flow of urine from the bladder into the ureters and potentially the kidneys, most commonly due to incompetent ureterovesical valve. Predisposes to pyelonephritis and renal scarring. More common in children, presenting with recurrent febrile UTIs. Diagnosis: voiding cystourethrogram (VCUG). Grading I-V. Management: observation with prophylactic antibiotics for low-grade, ureteral reimplantation or endoscopic injection for high-grade.

\medterm{Xanthogranulomatous Pyelonephritis}-- Chronic destructive renal infection with replacement of parenchyma by lipid-laden macrophages. Associated with obstructing stones and Proteus infection. Enlarged kidney with scattered low-attenuation areas on CT resembling a bear paw. Treated with nephrectomy as the kidney is non-functional and harbors chronic infection.
